problem:  
Let \( n \ge 3 \) and let \( S^{n+1} \subset \mathbb{R}^{n+2} \) carry the round metric.  
Fix a closed subgroup \( G \subset O(n+2) \) acting isometrically on \(S^{n+1}\) with cohomogeneity \(>1\).  
Denote by \(\mathcal{M}_G(S^{n+1})\) the set of embedded minimal hyperspheres \(\Sigma \subset S^{n+1}\) that are \(G\)-invariant, i.e. each \(\Sigma\) is diffeomorphic to \(S^n\) and invariant under the given group action.

**Open Problem:**  
Does there exist a subgroup \( G \) with cohomogeneity \(>1\) for which  
\[
\mathcal{M}_G^{A_0} \;=\; \{\Sigma \in \mathcal{M}_G(S^{n+1})\mid \operatorname{Area}(\Sigma)\le A_0\}/\mathrm{Isom}(S^{n+1})
\]
is infinite for some finite \( A_0>0 \)?  
In other words, can there exist infinitely many *noncongruent*, \(G\)-invariant, embedded minimal hyperspheres of uniformly bounded area in the standard sphere?  
Equivalently, does the imposition of symmetry under a fixed compact group ever *prevent* or *allow* blow-up of the area spectrum among \(G\)-invariant minimal hyperspheres in \(S^{n+1}\)?  

Why is it a "good" problem:  
This problem addresses a profound and currently unresolved intersection of **symmetry, compactness, and multiplicity** in the theory of minimal hypersurfaces.  
Without symmetry constraints, the work of Marques–Neves and others shows that a generic closed manifold with positive Ricci curvature (including \(S^{n+1}\)) admits infinitely many embedded minimal hypersurfaces whose areas diverge to infinity.  
However, when one restricts to hypersurfaces invariant under a compact Lie group \(G\), the variational space collapses to functions or sections on a *lower-dimensional orbit space*. This reduction introduces strong analytic constraints—it might drastically alter index growth, sweepout families, or the spectrum of first variations.  
The condition of **cohomogeneity \(>1\)** ensures that the \(G\)-orbit space retains nontrivial geometry rather than being one-dimensional, thereby preventing degeneracy to the rotationally symmetric case (where classifications are easier and often finite).  

Proving either finiteness or infinitude here would extend Sharp’s compactness theorem (which requires bounded index) or equivariant min–max constructions (which generally yield sequences of minimal hypersurfaces with growing area). The problem asks whether symmetry itself can play the role of an “effective index bound,” or conversely whether symmetric min–max sequences can still proliferate within a fixed area window.  

Conceptually, this is a bridge between **Lie group representation geometry, geometric measure theory, and analytic compactness**. Its statement is elementary—“can symmetry force compactness for bounded-area minimal hyperspheres in \(S^{n+1}\)?”—yet its resolution likely demands new insights into the equivariant Almgren–Pitts framework, Morse index bounds under symmetry, and the topology of orbit spaces.  
It exemplifies the ideal of a good mathematical problem: clear, simple, and accessible in formulation, but opening deep and interconnected fronts across geometry, analysis, and topology.
